\section{Actual Chern forms of model vector bundles}
\label{r10c:section}

Let $K$ be a complete nontrivially discretely valued field, let
$R=K^\circ$, and let $\mathcal X$ be an integral projective flat
$R$-scheme with smooth projective generic fiber $X$ of pure dimension $n$.
All analytic forms below are Mihatsch actual delta-forms on the ordinary
topology of $X^{\rm an}$; their sheaf is denoted $\mathcal B$ and
$D=d'd''$.  In the intrinsic trigrading, a form in
$\mathcal B^{0,0,r}$ has bidegree $(r,r)$.
The special fiber and its components need not be smooth, reduced, or
normal.  All Chow operations are rational and operational unless
otherwise stated.

\subsection{A scalar-density detector for pure delta forms}

\begin{lemma}[Detection with model Hessians]
\label{r10c:detector}
Let $\Omega\in\mathcal B^{0,0,r}(X^{\rm an})$ satisfy
$d'\Omega=d''\Omega=0$, with $0\leq r\leq n$, and put $q=n-r$.  Suppose that for every
tuple of global rational model functions $\phi_1,\ldots,\phi_q$,
\[
                 \Omega\wedge D\phi_1\wedge\cdots\wedge D\phi_q=0
\]
as a signed Radon measure.  Then $\Omega=0$ as an actual form.
For $q=0$, the hypothesis is the vanishing of the measure of $\Omega$.
\end{lemma}

\begin{proof}
An actual top-degree form has a locally finite signed measure: in finitely
many polyhedral charts over a compact set this follows by integration of
its smooth coefficient densities.  A partition of unity gives the same
measure as actual integration.  If that measure is zero, Mihatsch
Proposition 4.11 implies that the top form is zero.  This deals with $q=0$.

Assume that vanishing of the displayed measure has already been proved
with $k$ model-function slots replaced by arbitrary global original-smooth
functions.  Fix those $k$ functions and all but one of the remaining
model functions, and denote their product with $\Omega$ by $T$.
Then $T$ has bidegree $(n-1,n-1)$ and is separately closed.  For every
model function $\phi$ and every global original-smooth scalar $\psi$,
two applications of actual Stokes give
\[
 \int_{X^{\rm an}}\phi\,T\wedge D\psi
     =\int_{X^{\rm an}}\psi\,T\wedge D\phi=0.
\]
For fixed $T,\psi$, the left-hand side tests the single fixed finite
signed measure $T\wedge D\psi$ against $\phi$.  Scalar uniform density
therefore implies $T\wedge D\psi=0$.  Repeating this argument proves
\[
          \Omega\wedge D\chi_1\wedge\cdots\wedge D\chi_q=0
\tag{D}\label{r10c:smooth-hessians}
\]
for arbitrary global original-smooth functions $\chi_i$.
The density step concerns the continuous function \(\phi\) against one fixed signed measure, and consequently requires no convergence of Hessians.

If $\Omega$ is nonzero, choose an $n$-dimensional witnessing skeleton
in a moment chart, using Mihatsch Definition 3.8 and Lemma 4.4.
On an open polyhedron $\sigma$ of dimension $q$ its pure tridegree
means that the nonzero coefficient is scalar.  Refine the chart and
subdivide so that its moment coordinates $u_a=-\log|h_a|$ have
injective tangent map on $\sigma$; this follows from the finite-fiber
property of skeleton tropicalization.  At a point in the relative
interior of $\sigma$, choose a smooth cutoff equal to one nearby and
supported in the moment-chart open set.  It globalizes
\[
                         \chi=\tfrac12\sum_a u_a^2
\]
to an original-smooth scalar on $X^{\rm an}$.  Near the chosen point,
$D\chi=\sum_a d'u_a\wedge d''u_a$ is positive definite on the tangent
space of $\sigma$.  Thus $(D\chi)^q$ is a nonzero volume form there.
Coefficientwise wedge with smooth forms now shows that
$\Omega\wedge(D\chi)^q$ is a nonzero actual top form.  This contradicts
\eqref{r10c:smooth-hessians} and Proposition 4.11.
\end{proof}

\subsection{From operational numerical triviality to actual zero curvature}

For clarity, the CLD top measure agrees with the measure of the actual
wedge used above.  Locally let $u$ be a PWS model metric weight and let
$T$ be a product of previously treated actual model curvatures.  Such a
current is locally of order zero by its polyhedral description, so
coefficientwise multiplication by the continuous function $u$ is defined
and $[uT]=u[T]$. This equality is tested against original-smooth forms
and follows directly from the same order-zero polyhedral integrals
on both sides. Separate closedness and Mihatsch Proposition 4.11 give
\[
                  [Du\wedge T]=D[uT]=D\bigl(u[T]\bigr).
\]
Induction first identifies the actual wedge with the CLD recursive
product for semipositive model factors. For signed model metrics,
express them as differences of psh model
metrics and use multilinearity.  This establishes the comparison over
$K$ itself, before any scalar extension.

\begin{theorem}[Numerical model curvature vanishing]
\label{r10c:numerical-zero}
Let $P$ be a homogeneous polynomial of degree $r$ in the first Chern
classes of model line bundles $\mathcal L_1,\ldots,\mathcal L_t$ on
$\mathcal X$.  Suppose that
\[
       \xi=P\bigl(c_1(\mathcal L_{1,s}),\ldots,
                    c_1(\mathcal L_{t,s})\bigr)\in A^r(\mathcal X_s)_{\mathbb Q}
\]
has degree zero against every integral $r$-dimensional closed subvariety
of $\mathcal X_s$.  Then
\[
        P\bigl(c_1(\overline L_1),\ldots,c_1(\overline L_t)\bigr)=0
                 \quad\hbox{in }\mathcal B^{0,0,r}(X^{\rm an}).
\]
\end{theorem}

\begin{proof}
Let $\Omega$ be the displayed actual model polynomial.  Given $q=n-r$
model functions, choose a common proper model refinement
$f:\mathcal Y\to\mathcal X$ on which they are represented, after
clearing denominators, by vertical Cartier divisors $D_1,\ldots,D_q$.
One may work with formal completions; algebraic models and projective
common refinements suffice here.  The induced metrics of the original
bundles are unchanged by pullback.

For a point $x$ reducing to the generic point of a special component
$V\subset\mathcal Y_s$, the top wedge
$\Omega\wedge\prod_jD\phi_j$ has mass equal to the native CLD weight
$w_{\mathcal Y}(x)$ times
\[
 \deg\left(f_s^*\xi\cap
       \prod_{j=1}^{q}c_1(\mathcal O(D_j)|_V)\cap[V]\right).
\]
The proper operational projection formula identifies this degree with
\[
 \deg\left(\xi\cap f_{s*}
       \left(\prod_{j=1}^{q}c_1(\mathcal O(D_j)|_V)\cap[V]\right)\right)=0.
\]
The pushed-forward class has dimension $r$.  Terms whose images have
smaller dimension have zero proper pushforward; all remaining terms are
annihilated by the numerical hypothesis.  Operational pullback of a
Cartier polynomial is available for arbitrary $f_s$: neither flatness of
$f_s$ nor smoothness of $V$ has been used.  The CLD support and mass
formula now gives a zero measure, with all native multiplicities and
defect factors retained.  Lemma~\ref{r10c:detector} concludes the proof.
\end{proof}

\subsection{Descent along locally split morphisms}

\begin{lemma}[Descent for a locally split morphism]
\label{r10c:split-descent}
Let $a:Y\to X$ be a morphism of pure-dimensional boundaryless analytic
spaces.  Suppose $X$ has an ordinary open cover $(U_i)$ and morphisms
$s_i:U_i\to Y$ with $a\circ s_i=\operatorname{id}_{U_i}$.
For $\Theta\in\mathcal B(Y)$ the following are equivalent:
\begin{enumerate}
\item There is an actual form $\theta$ on $X$ with $a^*\theta=\Theta$.
\item On $Y\times_XY$, $p_1^*\Theta=p_2^*\Theta$.
\end{enumerate}
When the fiber product is pure-dimensional and the displayed pullbacks
are defined, the descended form is unique and is given by
$\theta|_{U_i}=s_i^*\Theta$.
\end{lemma}

\begin{proof}
The forward implication is functoriality.  Conversely pull the equality
back along $(s_i,s_j):U_i\cap U_j\to Y\times_XY$ to obtain equality
of the two local forms.  They glue by the sheaf property.  Pulling the
same equality along
$(\operatorname{id}_Y,s_i\circ a):a^{-1}(U_i)\to Y\times_XY$ proves
$\Theta=a^*s_i^*\Theta$ on $a^{-1}(U_i)$, hence $\Theta=a^*\theta$.
The local sections also prove uniqueness and injectivity of $a^*$ on
actual forms.  
\end{proof}

\begin{theorem}[Actual model Chern forms by local flag sections]
\label{r10c:chern}
Let $\mathcal E_1,\ldots,\mathcal E_s$ be vector bundles on
$\mathcal X$ and $P$ a homogeneous rational polynomial of degree $r$
in their Chern classes.  There is a unique actual form
\[
              c_{\mathcal X}(P)\in\mathcal B^{0,0,r}(X^{\rm an})
\]
whose pullback to the simultaneous full flag bundle equals the
corresponding symmetric polynomial in the model first Chern forms of
the tautological line quotients.  The form is separately closed,
independent of the chosen splitting presentation, and functorial under
morphisms of projective flat models with smooth projective generic
fibers.  Universal operational identities, in particular Whitney and
the additivity of the Chern character in exact sequences of model
bundles, hold as strict identities of actual forms.

If the special-fiber operational class $P(c(\mathcal E_{j,s}))$ is
numerically zero against every $r$-cycle, then $c_{\mathcal X}(P)=0$.
\end{theorem}

\begin{proof}
Put $a:\mathcal F=\prod_{\mathcal X}\operatorname{Flag}(\mathcal E_j)
\to\mathcal X$ and denote its generic fiber by $a:F\to X$.
Let $\Omega_F$ be the symmetric model-root polynomial.  On
$\mathcal F\times_{\mathcal X}\mathcal F$, the two root polynomials
represent the same operational Chow class by the two tautological
filtrations and Whitney's identity.  Their difference is therefore
operationally zero on the special fiber, before taking a degree or
capping with the fundamental class.  Theorem~\ref{r10c:numerical-zero},
applied to the smooth projective generic fiber $F\times_XF$, gives
\[
                         p_1^*\Omega_F=p_2^*\Omega_F
\]
as actual forms.

Simultaneously trivialize the finitely many generic bundles $E_j$ on
a Zariski open cover of $X$.  Their standard full flags give algebraic
sections $s_i:U_i\to F$, hence analytic sections on ordinary opens
$U_i^{\rm an}$.  These sections need not extend to all of the model
and need not preserve a model trivialization.  This is harmless:
they pull back an already defined global actual form, and Mihatsch
Theorem 4.5 applies to arbitrary analytic morphisms.  Applying
Lemma~\ref{r10c:split-descent} defines
$c_{\mathcal X}(P)|_{U_i^{\rm an}}=s_i^*\Omega_F$ and proves its
characterizing pullback equality and uniqueness.  Pullback preserves
the intrinsic trigrading and the differential operators, which proves
purity and separate closedness.

For any operational identity among finitely many model bundles, use
their simultaneous flag bundle.  The difference of the root polynomials
is operationally zero, so the numerical detector kills it upstairs.
Faithfulness supplied by local sections gives the identity downstairs.
The same argument on the product of two flag presentations proves
independence of presentation.  Functoriality follows by pulling the
model flags back and then using their characteristic equality and the
faithfulness of the new flag pullback.

Finally suppose $\xi=P(c(\mathcal E_{j,s}))$ is numerically zero.
For every integral $r$-cycle $W$ on $\mathcal F_s$, the operational
projection formula gives
\[
                 \deg(a_s^*\xi\cap[W])=
                 \deg(\xi\cap a_{s*}[W])=0.
\]
If the image drops dimension then $a_{s*}[W]=0$; otherwise the numerical
hypothesis applies.  Thus the root polynomial is numerically zero on
the whole flag special fiber.  Theorem~\ref{r10c:numerical-zero}
gives $\Omega_F=0$, and local sections give $c_{\mathcal X}(P)=0$.
Although $a_s$ happens to be smooth and flat here, the projection
argument is the operational one and does not use a smooth special fiber.
\end{proof}

The construction gives a homomorphism
\[
 K_0(\mathcal X)_{\mathbb Q}\longrightarrow
      \bigoplus_{r=0}^{n}\mathcal B^{0,0,r}(X^{\rm an}),
                  \qquad [\mathcal E]\longmapsto
                       \sum_r c_{\mathcal X}(\operatorname{ch}_r\mathcal E).
\]
Each degree factors through its special-fiber operational numerical
class.  This statement concerns classes of model bundles.  The dependence on the model is controlled by the numerical class of its special fiber.
